% Generado por herramientas/consolidar_a6.py. No editar a mano.
  \multicolumn{4}{@{}l}{\intersemibold Subdocumento 1 \textemdash{} Presentación de la empresa} \\
  Introducción & Claude / Gemini; OpenAI Codex & Síntesis del capítulo y conexión con el Formulario T-6 y los Subdocumentos 4, 7 y 12 & Guillermo Castillo: remisiones a la sección 1.5, al T-6 y al Art. 34, y a los Subdocumentos 4, 7 y 12, frente a su contenido vigente. \\
  1.1 & Claude / Gemini; OpenAI Codex & Trayectoria, inventario y puestos corporativos de atención & Alex Aravena: trayectoria, líneas de negocio, catálogo y antecedentes financieros; Guillermo Castillo: suma de 121 equipos y reparto de los 16 puestos del NOC y la mesa. \\
  1.2 & Claude / Gemini; OpenAI Codex & Organigrama, antigüedad, turnos y respaldos del NOC & Leandro Chamorro: organigrama, dotación por área y sede, antigüedad y turnos del NOC. \\
  1.3 & Claude / Gemini; OpenAI Codex & Redacción de gobierno y responsabilidades & Guillermo Castillo: periodicidad mensual del Comité de Calidad y Procesos frente al SD6, sección 6.1.6; coberturas de 80~\% global, 70~\% de lógica de negocio (RT-04.11) y 80~\% del código modificado frente al SD6 y al T-10; auditorías internas trimestrales; y pruebas de intrusión independientes frente al RT-11.20. \\
  1.4 & Claude / Gemini; OpenAI Codex & Comparación de experiencia y SLA & Patricio Henríquez: certificaciones, experiencia, SLA y conocimiento normativo del caso. \\
  1.5 & Claude / Gemini; Claude Code; OpenAI Codex & Organización por ámbitos y coordinación de roles & Alex Aravena: ámbitos y roles del proyecto frente al T-15. \\
  1.6 & Claude / Gemini; OpenAI Codex & Redacción de alcance de alianzas & Álvaro Catalán: alianzas, vigencias y plan de renovación. \\
  T-6 & Claude / Gemini; OpenAI Codex & Formato, introducción y referencias & Tomás Pérez: los once campos de los tres proyectos y su coherencia con la sección 1.4. \\
  \midrule
  \multicolumn{4}{@{}l}{\intersemibold Subdocumento 2 \textemdash{} Comprensión del problema y de la necesidad} \\
  Introducción & Claude / Gemini / Codex & Síntesis del diagnóstico y conexión con el alcance y los anexos & Alex Aravena: revisó la presentación del problema y su conexión con el alcance y los anexos. \\
  2.1 Resumen & Claude / Gemini / Codex & Síntesis y contraste de cifras y sistemas & Alex Aravena y Leandro Chamorro: revisión del resumen y de sus cifras. \\
  2.2 Problema & Claude / Gemini / Codex & Diagnóstico, tensiones y modelos de atención & Alex Aravena y Leandro Chamorro: revisión de problemas y respaldo normativo. \\
  2.3 Procesos & Claude / Gemini / Codex & Flujos AS-IS y corrección de figuras & Alex Aravena y Leandro Chamorro: revisión de flujos y cálculos. \\
  2.4 Actores & Claude / Gemini & Disposición tabular de actores & Alex Aravena y Leandro Chamorro: validación de los actores frente al catálogo. \\
  2.5 Síntesis & Claude / Gemini / Codex & Requerimientos y doce restricciones & Alex Aravena y Leandro Chamorro: revisión de supuestos y restricciones. \\
  Anexo~2.1 & Claude / Gemini & Maquetación del listado de requerimientos & Alex Aravena y Leandro Chamorro: validación del listado frente al catálogo. \\
  Anexo~2.2 & Claude / Gemini / Codex & Supuestos, restricciones y ajuste de S-09 & Alex Aravena y Leandro Chamorro: contraste con caps. 10, 11 y 16.1; revisión de S-09. \\
  Anexo~2.3 & Claude / Gemini & Maquetación de actores y sistemas legados & Alex Aravena y Leandro Chamorro: revisión de actores, cifras y sistemas. \\
  \midrule
  \multicolumn{4}{@{}l}{\intersemibold Subdocumento 3 \textemdash{} Esquema de solución y alcance} \\
  3.1 & Codex, Claude Code & Contraste con las Bases y redacción. Incorporación de la cadena de frío en ruta. & Patricio Henríquez: cotejó los tres problemas del caso con el Subdocumento 2 y los hitos de ambas etapas con el T-18. \\
  3.2 & Codex, Claude Code & Contraste con las Bases y los Subdocumentos 1 y 2. Redacción y propuestas para las reglas RNG-04 y RNG-15. Incorporación de la cadena de frío en ruta. Coherencia del alcance con el Formulario T-12. & Patricio Henríquez: contrastó exclusiones y restricciones con el Caso Logística y cotejó los plazos y el corte de las 14:00 con S-21 y RNG-15. \\
  3.3 & Codex, Claude Code & Descripción de las figuras, actualización del esquema de trazabilidad del lote y revisión de la continuidad. Incorporación de la cadena de frío en ruta. & Patricio Henríquez: cotejó el recorrido del lote, la decisión de Calidad y la continuidad durante un corte con el Anexo 3.G y el Subdocumento 4. \\
  3.4 & Codex, Claude Code & Redacción del desarrollo, despliegue y operación. Revisión de reversión, actores, módulos y referencias. Actualización de la vista general y elaboración del esquema de emplazamiento. Incorporación de la cadena de frío en ruta. & Patricio Henríquez: cotejó módulos y emplazamiento con los anexos 4-D y 4-N, el alta y reemplazo de conductores con el Subdocumento 4 y la dotación con el Anexo 4-W.7 y el T-15. \\
  Anexo 3.A & Codex, Claude Code & Conversión a tablas del catálogo funcional. Incorporación de la cadena de frío en ruta. Revisión de resultados y precondiciones del catálogo. & Guillermo Castillo: cotejó la columna Origen con los apartados que respaldan los requerimientos funcionales. \\
  Anexo 3.B & Codex, Claude Code & Conversión a tablas del catálogo no funcional. Incorporación de la cadena de frío en ruta. & Guillermo Castillo: cotejó los umbrales con sus fuentes y comprobó los métodos de verificación. \\
  Anexo 3.C & Codex, Claude Code & Redacción de supuestos. Incorporación de la cadena de frío en ruta. & Guillermo Castillo: comprobó antecedente, impacto y validación de los supuestos y cotejó S-21 con D-03 y RNG-15. \\
  Anexo 3.D & Codex, Claude Code & Redacción del registro de exclusiones. Precisión de la responsabilidad del CLIENTE en la compra del kit de cadena de frío de los camiones. & Guillermo Castillo: cotejó las exclusiones con el capítulo 11 del Caso Logística y comprobó su efecto sobre el alcance. \\
  Anexo 3.E & Codex, Claude Code & Redacción del registro de restricciones operacionales. & Guillermo Castillo: contrastó las restricciones con el Caso Logística y las Bases y comprobó su efecto sobre el diseño o la operación. \\
  Anexo 3.F & Codex, Claude Code & Redacción del registro de decisiones de alcance y diseño. & Guillermo Castillo: contrastó las decisiones de alcance con los antecedentes citados y comprobó su fundamento. \\
  Anexo 3.G & Codex, Claude Code & Redacción de reglas y propuestas para RNG-04 y RNG-15. Incorporación de la cadena de frío en ruta. & Guillermo Castillo: comprobó la autoridad de Calidad en RNG-04 y cotejó la promesa de RNG-15 con S-21. \\
  Anexo 3.H & Codex, Claude Code & Redacción del registro de vacíos y consultas. & Guillermo Castillo: cotejó los antecedentes de las consultas y comprobó su efecto sobre el alcance y su responsable. \\
  Anexo 3.I & Codex, Claude Code & Redacción de la estrategia de participación y del catálogo de actores del sistema, con sus interfaces e interacciones. & Guillermo Castillo: contrastó los actores con el Subdocumento 2 y distinguió la participación en el proyecto del acceso a aplicaciones. \\
  Anexo 3.J & Codex, Claude Code & Redacción de criterios de aceptación. Incorporación de la cadena de frío en ruta. & Guillermo Castillo: cotejó el umbral de OTIF del mes 32 con el cuerpo del Subdocumento 3 y la fila R18-04 del T-17. \\
  Anexo 3.K & Codex, Claude Code & Redacción del glosario de siglas y códigos. & Guillermo Castillo: cotejó las definiciones del glosario con el uso de los códigos en el cuerpo y los anexos. \\
  T-12 & Codex, Claude Code & Generación de la matriz de cumplimiento. Incorporación de la cadena de frío en ruta. Revisión de componentes, paquetes, referencias y estado de cumplimiento por requisito. & Patricio Henríquez: cotejó componentes con el Subdocumento 4, paquetes con el T-14 y pruebas y criterios con la matriz del T-13. \\
  \midrule
  \multicolumn{4}{@{}l}{\intersemibold Subdocumento 4 \textemdash{} Arquitectura lógica y física de la solución} \\
  Introducción & OpenAI Codex & Síntesis de la arquitectura y conexión con anexos y Formulario T-11. & Alex Aravena: revisó la síntesis de la arquitectura y su conexión con los anexos y el Formulario T-11. \\
  4.1 & OpenAI Codex & Corrección de identidad, trazabilidad, síntesis y variantes vectoriales de figuras. & Patricio Henríquez: trazabilidad de los requisitos con el T-12 y nombres de módulos en las figuras. \\
  4.2 & OpenAI Codex & Redacción y organización de la arquitectura física y sus diagramas. & Álvaro Catalán: infraestructura, conexiones, despliegue y vigencia de la credencial en reparto. \\
  4.3 & OpenAI Codex; Claude & Redacción y organización de los centros de datos y su recuperación; asistencia en el plano del recinto técnico. & Álvaro Catalán: revisó la descripción de los centros de datos y las medidas de recuperación. \\
  4-A & OpenAI Codex & Organización y descripción de los eventos. & Álvaro Catalán: revisó la descripción de los eventos y su relación con la operación. \\
  4-B & OpenAI Codex & Redacción de las reglas de integración y respuesta ante fallas. & Álvaro Catalán: revisó las reglas de integración y la respuesta ante fallas. \\
  4-C & OpenAI Codex & Organización de los escenarios de carga. & Álvaro Catalán: revisó la descripción de los escenarios de carga. \\
  4-D & OpenAI Codex & Consolidación del catálogo funcional y síntesis por contexto. & Patricio Henríquez: catálogo funcional y nombres de módulos frente al SD3. \\
  4-E & OpenAI Codex & Matriz de requisitos críticos individuales y criterios de aceptación. & Tomás Pérez: requisitos críticos, pruebas y criterios de aceptación frente al T-13. \\
  4-F & OpenAI Codex & Descripción de los límites entre módulos. & Álvaro Catalán: revisó la claridad de los límites y las relaciones entre módulos. \\
  4-G & OpenAI Codex & Coherencia del contrato de captura con la ruta declarada. & Alex Aravena: contrato de captura frente a la ruta declarada. \\
  4-H & OpenAI Codex & Coherencia del horario de integración con la ruta declarada. & Alex Aravena: horario de integración frente a la ruta declarada. \\
  4-I & OpenAI Codex & Organización de los volúmenes de información. & Álvaro Catalán: revisó la presentación de los volúmenes y sus cálculos. \\
  4-J & OpenAI Codex & Descripción de la operación sin conexión. & Álvaro Catalán: revisó las funciones disponibles y sus límites durante un corte. \\
  4-K & OpenAI Codex & Redacción de las reglas de reconciliación. & Álvaro Catalán: revisó la descripción de las reglas para conciliar la información. \\
  4-L & OpenAI Codex & Organización de las decisiones solicitadas por el caso. & Álvaro Catalán: revisó la correspondencia de las decisiones con lo solicitado en el caso. \\
  4-M & OpenAI Codex & Protocolos de manifiesto, reloj, revocación y ruta de 14,5 h. & Guillermo Castillo: vigencia del manifiesto, reloj, revocación y relevos en la ruta de 14,5 h. \\
  4-N & OpenAI Codex & Consolidación del detalle de interacción y nombres de módulos. & Patricio Henríquez: detalle de interacción y nombres de módulos frente al SD3. \\
  4-O & OpenAI Codex & Consistencia de ADR-06 y ADR-23 con identidad y backend. & Alex Aravena: ADR-06 y ADR-23 frente a la identidad y el backend. \\
  4-P & OpenAI Codex & Verificación documental de existencia, soporte y compatibilidad lógica. & Guillermo Castillo: lanzamiento y soporte de cada producto contra sus fuentes oficiales. \\
  4-Q & OpenAI Codex & Organización y descripción de las amenazas. & Álvaro Catalán: revisó la relación entre las amenazas y los componentes de la solución. \\
  4-R & OpenAI Codex & Consolidación de custodia y evidencia de controles lógicos. & Álvaro Catalán: custodia y evidencia de los controles lógicos. \\
  4-S & OpenAI Codex & Organización de las vistas y sus relaciones. & Álvaro Catalán: revisó la coherencia entre las vistas y sus relaciones. \\
  4-T & OpenAI Codex & Descripción del desempeño y sus criterios de aceptación. & Álvaro Catalán: revisó la presentación de los escenarios y criterios de desempeño. \\
  4-U & OpenAI Codex & Organización de los documentos y las evidencias de entrega. & Álvaro Catalán: revisó la relación entre los documentos y las evidencias de entrega. \\
  4-V & OpenAI Codex & Especificación de pruebas de fronteras, frío y auditoría. & Tomás Pérez: pruebas de fronteras, frío y auditoría frente al SD9. \\
  4-W & OpenAI Codex & Organización de los cálculos de capacidad. & Álvaro Catalán: revisó la presentación de los cálculos y su relación con la capacidad propuesta. \\
  4-X & OpenAI Codex & Reorganización de detalle existente y traslado de las Figuras 19--23, sin alterar sus imágenes. & Álvaro Catalán: revisó la presentación de las figuras y su relación con los ambientes de despliegue. \\
  T-11 & OpenAI Codex & Organización y revisión del inventario de equipamiento. & Álvaro Catalán: revisó el inventario y su correspondencia con la solución. \\
  \midrule
  \multicolumn{4}{@{}l}{\intersemibold Subdocumento 5 \textemdash{} Modelo y gestión de datos} \\
  5.1 & GPT; Codex & Modelo conceptual y lógico; trazabilidad y modelos por dominio; figuras 5.1–5.10. Incorporación de la historia térmica por SSCC. & Álvaro Catalán y Leandro Chamorro. Revisión de los modelos propuestos. Verificación de figuras acorde a los modelos propuestos. \\
  5.2 & GPT; Codex & Persistencia, CAP, integridad, calidad, explotación, retención e intercambio. Incorporación de la historia térmica por SSCC. & Maximiliano Miño. Revisíon de los motores, la autoridad de escritura y las reglas de integridad, calidad, retención e intercambio. \\
  5.3 & GPT; Codex & Volúmenes, saneamiento, herramientas, ensayos, conciliación y corte de migración. & Leando Chamorro y Maximiliano Miño: Revisión de volúmenes, las transformaciones y la secuencia de migración, incluidos los criterios de conciliación, corte y reversión. \\
  5.4 & GPT; Codex & Escenarios de carga, índices, particiones, cachés, concurrencia y pruebas. Incorporación de la historia térmica por SSCC. & Álvaro Catalán y Maximiliano Miño: Revisión cálculos de carga, consultas críticas y estrategias de índices, particionado, caché y concurrencia. \\
  5-A & GPT; Codex & Diccionario de atributos, políticas, identidades y campos de integración. Incorporación de la historia térmica por SSCC. Recuento del diccionario. & Leandro Chamorro. Revisión definiciones, tipos, claves y nulabilidad de los atributos y su coherencia con el modelo. \\
  5-B & GPT & Cardinalidades, módulos e integridad referencial. & Álvaro Catalán. Revisión cardinalidades, la asignación por módulo y las restricciones de integridad entre entidades. \\
  5-C & GPT; Codex & Dominios de valores, validaciones y causas de rechazo. Incorporación de la historia térmica por SSCC. & Álvaro Catalán y Leandro Chamorro. Revisión los valores permitidos, las reglas de validación y las causas de rechazo de los datos. \\
  5-D & GPT & Retención, archivado y eliminación por dominio y atributo. & Álvaro Catalán y Maximiliano Miño. Revisión plazos de conservación y las condiciones de archivado y eliminación por dominio. \\
  5-E & GPT & Clasificación por sensibilidad y controles de protección. & Maximiliano Miño y Leandro Chamorro. Revisión clasificación de los datos y su correspondencia con los controles de acceso y protección. \\
  5-F & GPT; Codex & Índices, columnas, restricciones y particionado. Incorporación de la historia térmica por SSCC. & Leandro Chamorro. Revisión columnas indexadas, las restricciones y el particionado respecto de las consultas previstas. \\
  5-G & GPT & Claves de caché, vigencias e invalidación. & Álvaro Catalán. Revisión claves, los tiempos de vigencia y las condiciones de invalidación y reconstrucción de las cachés. \\
  5-H & GPT & Fórmulas de indicadores, linaje y catálogo por fase. & Maximiliano Miño. Revisión fórmulas, las unidades y el origen de los datos utilizados en cada indicador. \\
  5-I & GPT; Codex & Trazabilidad de RT-05 y evidencia documental. Correspondencia de RT-05.01 con el diccionario. & Leandro Chamorro y Álvaro Catalán. Revisión correspondencia de los requisitos RT-05 con las secciones y evidencias citadas. \\
  5-J & GPT & Protocolos de aceptación de datos y pruebas. & Álvaro Catalán y Maximiliano Miño. Revisión escenarios propuestos, sus criterios de aceptación y las evidencias previstas para las pruebas. \\
  5-K & GPT & Mapeo de origen y destino, olas y ventana de corte. & Leandro Chamorro y Maximiliano Miño. Revisión mapeo de datos, los cálculos por ola y la secuencia de la ventana de corte. \\
  5-L & GPT & Correspondencia del diccionario con los modelos del cuerpo. & Maximiliano Miño. Revisión correspondencia de los modelos complementarios con las figuras del cuerpo y el diccionario. \\
  \midrule
  \multicolumn{4}{@{}l}{\intersemibold Subdocumento 6 \textemdash{} Metodologías} \\
  Introducción & Chat GPT; Claude Code & Redacción de la introducción y conexión con SD2 a SD4 y SD7 a SD9 & Tomás Pérez: que cada remisión a otro subdocumento corresponda a su contenido vigente \\
  6.1 Metodología de Gestión de Proyectos & Chat GPT; Codex; Claude Code & Adaptación del PMBOK, control de cambios del Art. 72°, escalamiento, nivelación y Figuras 6.1 y 6.2 & Maximiliano Miño: frontera fijo/adaptable, plazos del control de cambios y disparadores de escalamiento frente al SD8; Guillermo Castillo: coherencia de los comités con el Art. 71° y el T-14 \\
  6.2 Metodología de Desarrollo Software & Chat GPT; Codex; Claude Code & Comparación de enfoques, gestión de requerimientos, deuda técnica, condiciones de Puelche y figura de la cadena de entrega & Tomás Pérez: paquetes y pruebas citados en la tabla de condiciones propias del caso frente al T-14 y el T-13; Maximiliano Miño: puertas y ambientes frente al SD4 y el SD9 \\
  Anexo 6.A — Interesados del proyecto & Chat GPT; Codex; Claude Code & Registro de los 19 interesados & Guillermo Castillo y Maximiliano Miño: niveles de influencia e interés frente al SD2 y estrategia frente al SD3 \\
  Anexo 6.B — Adquisiciones asociadas al proyecto & Chat GPT; Codex; Claude Code & Registro de adquisiciones & Guillermo Castillo y Maximiliano Miño: responsables, plazos y evidencias de las adquisiciones \\
  Formulario T-9 & Chat GPT; Claude Code & Procedimientos, umbrales, matrices y plantillas de gestión & Maximiliano Miño: umbrales de valor ganado y plazos del control de cambios; Tomás Pérez: matriz de comunicaciones frente al Anexo 6.A \\
  Formulario T-10 & Chat GPT; Claude Code & Ambientes, revisiones, deuda técnica y artefactos del desarrollo & Tomás Pérez: ambientes y aprobaciones frente al SD4; Guillermo Castillo: controles y entregas de cada fase frente a la programación del T-15 \\
  \midrule
  \multicolumn{4}{@{}l}{\intersemibold Subdocumento 7 \textemdash{} Plan de trabajo, EDT, cronograma e implantación} \\
  Introducción & Claude Code & Síntesis de las Bases y los capítulos de la oferta & Leandro Chamorro: revisó que lo presentado coincidiera con lo mostrado en el subdocumento y se alineara con lo solicitado en las Bases. \\
  7.1 EDT & Claude Code; Codex & Redacción, conteos y ajuste de las figuras de la EDT & Leandro Chamorro: comprobó que la organización entre ramas y paquetes de trabajo estuviera bien organizada y fuera coherente. \\
  7.2 Plan de trabajo & Claude Code; Codex & Secuenciamiento, dependencias y frentes de trabajo & Leandro Chamorro: revisó que el orden de las actividades concordara con lo establecido en la EDT y se mantuviera consistente con los demás subdocumentos. \\
  7.3 Cronograma e implantación & Claude Code; Codex & Calendario, Gantt, CPM/PERT y esquemas de implantación & Leandro Chamorro: comprobó que las metas coincidieran y se alinearan con lo solicitado en las Bases y con su correspondencia con los paquetes. \\
  Anexo 7.A & Claude Code; Codex & Cálculo de meses contractuales y ventanas de congelamiento según fecha de inicio & Tomás Pérez: revisó que lo mostrado en el anexo correspondiera y se alineara con lo entregado por la EDT y el cronograma. \\
  Anexo 7.B & Claude Code; Codex & Consolidación de la tabla de dependencias entre paquetes de trabajo & Tomás Pérez: revisó que las dependencias entre paquetes coincidieran lógicamente con lo mostrado en el calendario y los hitos. \\
  Anexo 7.C & Claude Code; Codex & Mapeo del momento de verificación de los resultados de aceptación del caso & Tomás Pérez: revisó que lo mostrado en el anexo correspondiera o se alineara con las actividades del cronograma, y que fueran correctos las metas, momentos y paquetes relacionados. \\
  Anexo 7.D & Claude Code; Codex & Ubicación de las cinco innovaciones en los paquetes de la EDT y el cronograma & Patricio Henríquez: revisó que la correspondencia entre cada innovación y los paquetes de la EDT, junto con sus hitos relacionados, coincidiera entre sí y consideró las actualizaciones realizadas en otros subdocumentos. \\
  Anexo 7.E & Claude Code; Codex & Trazabilidad de módulos, servicios comunes e interfaces hacia los paquetes de la EDT & Tomás Pérez: comprobó que los paquetes de trabajo tuvieran una relación directa con los módulos por desarrollar. \\
  Anexo 7.F & Claude Code; Codex & Supuestos de planificación, disparadores y enlace con la gestión de riesgos & Tomás Pérez: revisó que los supuestos estuvieran correctamente relacionados con sus respectivos paquetes y con los responsables de tratamiento. \\
  Formulario T-14 & Claude Code; Codex & Diccionario de la EDT, ventanas y cartas Gantt & Patricio Henríquez: comprobó que cada paquete estuviera definido de acuerdo con lo solicitado en las Bases y revisó cómo se conecta y relaciona con las ramas de la EDT. \\
  Formulario T-15 & Claude Code; Codex & HH, nivelación, actividades de 8 a 80 horas, PERT y frentes & Maximiliano Miño: revisó las estimaciones de HH y la nivelación de las actividades, y comprobó que los cálculos relacionados con la estimación PERT fueran correctos. \\
  Formulario T-18 & Claude Code; Codex & Plan de olas, marcha blanca, estabilización y reversión & Leandro Chamorro: contrastó las olas y los criterios relacionados con la marcha blanca y comprobó que el cronograma se alineara con ellos. \\
  \midrule
  \multicolumn{4}{@{}l}{\intersemibold Subdocumento 8 \textemdash{} Plan de riesgos} \\
  Introducción & Gentle-Shell y ChatGPT & Conexión con operación, alcance y formularios & Maximiliano Miño: revisión del alcance de la introducción y de su correspondencia con la operación y los entregables. \\
  8.1 Plan de riesgos & Gentle-Shell y ChatGPT & Roles, escalas, ciclo y cierre & Guillermo Castillo: roles, escalas de probabilidad e impacto, ciclo y cierre frente al SD6 y el T-16. \\
  8.2 Identificación y Análisis de Riesgos & Gentle-Shell y ChatGPT & RBS, exposición cualitativa y valor esperado en HH & Maximiliano Miño: enunciado de los riesgos, RBS y valoración de probabilidad, impacto y exposición; Patricio Henríquez: estimación de tres valores. \\
  8.3 Plan de Acción a Riesgos & Gentle-Shell y ChatGPT & Respuestas, contingencia y gestión & Alex Aravena: correspondencia entre la valoración, las respuestas, las contingencias y las reservas frente al SD7. \\
  Anexo 8.A & Gentle-Shell y ChatGPT & Proponer respuestas y estrategias para los riesgos, junto con un responsable para cada uno & Maximiliano Miño: revisión de la pertinencia de las estrategias, respuestas y responsables de cada riesgo. \\
  Anexo 8.B & Gentle-Shell y ChatGPT & Proponer y validar los cálculos relacionados con la estimación de tres valores & Patricio Henríquez: estimaciones de esfuerzo, cálculos PERT y valores esperados por riesgo. \\
  Anexo 8.C & Gentle-Shell y ChatGPT & Fundamentar la valoración inicial de impacto y su relación con el PERT y las dimensiones no monetarias & Guillermo Castillo: revisión de escalas y consistencia de los criterios de impacto. \\
  Anexo 8.D & Gentle-Shell y ChatGPT & Organizar dependencias operativas, responsables, continuidad y evidencias habilitantes & Maximiliano Miño: revisión de riesgos, dependencias y correspondencia con la operación y los entregables. \\
  Anexo 8.E & Gentle-Shell y ChatGPT & Definir controles técnicos, responsables, pruebas y condiciones de cierre & Alex Aravena: revisión de la correspondencia entre controles, respuestas y evidencia de cierre. \\
  Anexo 8.F & Gentle-Shell y ChatGPT & Desarrollar respuestas, estimaciones de esfuerzo y proyección del riesgo residual & Alex Aravena: revisión de respuestas, contingencias y reservas; Patricio Henríquez: revisión de estimaciones y cálculos. \\
  Anexo 8.G & Gentle-Shell y ChatGPT & Analizar modos de falla, efectos, prioridades y controles posteriores & Patricio Henríquez: revisión de estimaciones y cálculos; Guillermo Castillo: revisión de criterios de priorización. \\
  Anexo 8.H & Gentle-Shell y ChatGPT & Distribuir reservas, perfiles y esfuerzo en el tiempo, y explicitar supuestos y sensibilidad & Alex Aravena: revisión de la correspondencia entre reservas y esfuerzo frente al SD7. \\
  Formulario T-16 & Gentle-Shell y ChatGPT & Mejorar la redacción de las descripciones de riesgos y proponer valores de probabilidad e impacto, y estrategias o respuestas & Maximiliano Miño: enunciados, valoraciones de probabilidad e impacto y estrategias de respuesta; Tomás Pérez: contingencias y responsables. \\
  \midrule
  \multicolumn{4}{@{}l}{\intersemibold Subdocumento 9 \textemdash{} Plan de calidad} \\
  Introducción & Claude Code & Redacción a partir de las Bases (BA/BTT) y del Caso Logística & Patricio Henríquez y Guillermo Castillo. Verificamos la alineación de los antecedentes del Caso 02 (retiro sanitario y excursión térmica), la consistencia de las referencias cruzadas con SD1 a SD7 y la trazabilidad frente a los artículos de las BA/BTT. \\
  9.1 Plan de Calidad & Claude Code; Codex & Redacción, umbrales y actualización del costo de calidad & Patricio Henríquez y Guillermo Castillo. Revisamos los 4 objetivos de calidad, las 8 características de ISO/IEC 25010 y las 8 métricas de código bloqueantes de la Tabla 9.4. También comprobamos el cuadre de las 12.058,66 HH del costo de la calidad y la sintaxis lógica de las Figuras 9.1 y 9.2. \\
  9.2 Estrategia de Aseguramiento de Calidad & Claude Code; Codex & Reglas de calidad, actualización del catálogo, matriz y alcance por etapa & Patricio Henríquez y Guillermo Castillo. Verificamos las reglas y los bloqueos de las compuertas G0–G6, el dimensionamiento de carga con un factor de 1,5$\times$ y un peak de 21,99 TPS, y el catálogo de la revisión anterior. Además, revisamos el protocolo de datos sintéticos y anonimizados y la consistencia técnica de las Figuras 9.3 a 9.8. \\
  9.3 Alineación con Plan de Trabajo & Claude Code; Codex & Conciliación con T-14/T-15, calendario de campañas y actualización de diagramas & Patricio Henríquez y Guillermo Castillo. Comprobamos la conciliación de las 12.058,66 HH en los paquetes EDT del T-15 y revisamos el calendario de hitos H1–H12. También verificamos que las pruebas quedaran fuera de los períodos de congelamiento de los meses 8, 11 y 20, el cálculo de holguras ante subsanaciones y el diagrama Gantt de la Figura 9.9. \\
  Anexo 9.A & Claude Code; Codex & Redacción y organización de las métricas de calidad del producto y de los datos & Patricio Henríquez y Guillermo Castillo: revisaron las 42 métricas, sus umbrales, métodos de verificación y correspondencia con los paquetes de prueba. \\
  Anexo 9.B & Claude Code; Codex & Organización de las puertas de calidad y sus reglas de bloqueo & Patricio Henríquez y Guillermo Castillo: revisaron las puertas G0–G6, las herramientas por lenguaje y la evidencia exigida para cada promoción. \\
  Anexo 9.C & Claude Code; Codex & Cálculo y desglose del costo de la calidad en HH & Patricio Henríquez y Guillermo Castillo: revisaron los cálculos de prevención, evaluación y fallas, y la conciliación de las 12.058,66 HH de conformidad con el T-15, sin doble conteo. \\
  Anexo 9.D & Claude Code; Codex & Trazabilidad de los resultados del caso hasta sus pruebas e hitos & Patricio Henríquez y Guillermo Castillo: revisaron la correspondencia de los 16 resultados con las pruebas del T-13, los paquetes de medición y los hitos de aceptación, incluida la cadena de frío de R18-03. \\
  Anexo 9.E & Claude Code; Codex & Organización de las pruebas durante la construcción e integración & Patricio Henríquez y Guillermo Castillo: revisaron la programación de pruebas por actividad, las condiciones de integración en QA y promoción a Preproducción, y la evidencia vinculada al T-13, T-14 y T-15. \\
  Formulario T-13 & Claude Code; Codex & Calendario desde T-15 y trazabilidad de las validaciones de frío & Patricio Henríquez y Guillermo Castillo. Verificamos la estructura del plan de pruebas bajo ISO/IEC/IEEE 29119, los criterios de aceptación por tipo de prueba, los protocolos de resiliencia híbrida y las validaciones de frío. Asimismo, comprobamos la concordancia de las fechas y las HH con el formulario T-15. \\
  Formulario T-17 & Claude Code; Codex & Entregas, actas, holguras y protocolo de subsanación & Patricio Henríquez y Guillermo Castillo. Revisamos el protocolo de actas de aceptación técnica y administrativa, las condiciones de cierre de la marcha blanca establecidas en el Art. 17.3 y el esquema de aceptación diferida de los resultados de negocio. También contrastamos las entregas, actas y holguras de los diez hitos con entregable propio (todos salvo H6 y H11, que solo habilitan las marchas blancas) con el formulario T-15. \\
  \midrule
  \multicolumn{4}{@{}l}{\intersemibold Subdocumento 10 \textemdash{} Servicios de operación y niveles de servicio} \\
  Introducción & OpenAI Codex & Redacción de la introducción & Guillermo Castillo: síntesis del servicio y conexión con otros subdocumentos y anexos. \\
  10.1 & OpenAI Codex & Redacción, cálculo y diagramas & Guillermo Castillo: cobertura, horarios y dotación de mesa, NOC y SOC contra SD4 y T-15. \\
  10.2 & OpenAI Codex & Redacción y contraste de SLA & Guillermo Castillo: umbrales del Art. 78, relojes y escalamiento. \\
  10.A & OpenAI Codex & Memoria y reproducción de cálculo & Leandro Chamorro: entradas y resultados del cálculo de capacidad. \\
  10.B & OpenAI Codex & Indicadores y conservación & Leandro Chamorro: indicadores, fuentes de medición y retención. \\
  10.C & OpenAI Codex & Procedi\-mientos y coordinación & Guillermo Castillo: guías de recuperación, RTO y RPO contra SD4. \\
  10.D & OpenAI Codex & Contraste de dependencias & Guillermo Castillo y Leandro Chamorro: condiciones de habilitación. \\
  10.E & OpenAI Codex & Matriz de requisitos & Leandro Chamorro: localización de cada requisito. \\
  \midrule
  \multicolumn{4}{@{}l}{\intersemibold Subdocumento 11 \textemdash{} Planes en operación} \\
  11.1 & Codex & Texto, cálculos y figura de mantención & Bastián Trejo y Guillermo Castillo: revisión y validación de contenido, cálculos y figuras generadas. \\
  11.2 & Codex & Texto, cálculos y figura de operación & Bastián Trejo y Guillermo Castillo: revisión y validación de contenido, cálculos y figuras generadas. \\
  \midrule
  \multicolumn{4}{@{}l}{\intersemibold Subdocumento 12 \textemdash{} Equipo de trabajo, subcontrataciones y alianzas} \\
  Introducción & Claude Code; Codex & Articulación de los documentos de la oferta & Alex Aravena y Leandro Chamorro: correspondencia con los subdocumentos y formularios citados. \\
  12.1 & Claude Code; Codex & Síntesis de roles, dedicación y dotación & Alex Aravena y Leandro Chamorro: cotejo de las dedicaciones, las plazas y los diagramas con el Formulario T-15. \\
  12.2 & Claude Code; Codex & Distinción entre compras y servicios por resultado & Alex Aravena y Leandro Chamorro: contraste de paquetes, HH y responsables con la EDT. \\
  12.3 & Claude Code; Codex & Descripción del aporte de los fabricantes & Alex Aravena y Leandro Chamorro: comprobación del alcance y las vigencias de las alianzas. \\
  Anexo 12.A & Claude Code; Codex & Cruce de productos y roles & Alex Aravena y Leandro Chamorro: conciliación de los paquetes del T-14 con las intervenciones del T-15. \\
  Anexo 12.B & Claude Code; Codex & Organización de los antecedentes curriculares & Alex Aravena y Leandro Chamorro: cotejo de experiencia, certificaciones y vigencias por integrante. \\
  Anexo 12.C & Claude Code; Codex & Cálculo y presentación de dedicaciones & Alex Aravena y Leandro Chamorro: revisión de denominadores, plazas del perfil y períodos. \\
  Anexo 12.D & Claude Code; Codex & Redacción del compromiso individual & Alex Aravena y Leandro Chamorro: revisión de titulares, períodos y condiciones de reemplazo. \\
  Formulario T-8 & Claude Code; Codex & Consolidación de nominaciones y servicios externos & Alex Aravena y Leandro Chamorro: contraste de las dos declaraciones con el Subdocumento 12. \\
  \midrule
  \multicolumn{4}{@{}l}{\intersemibold Subdocumento 13 \textemdash{} Innovaciones} \\
  Introducción & Claude Code; OpenAI Codex & Síntesis, citas y mapa general & Guillermo Castillo: Se revisó la claridad de la presentación general y la comprensión del propósito de las cinco innovaciones. \\
  13.1 & Claude Code; OpenAI Codex & Contenido y recorrido INN-01 & Guillermo Castillo: Se examinó la redacción del seguimiento de vencimiento posentrega y la comprensión del problema y de la solución propuesta. \\
  13.2 & Claude Code; OpenAI Codex & Contenido y recorrido INN-02 & Guillermo Castillo: Se revisó la secuencia explicativa del recorrido de reproducción de incidentes de terreno y la claridad de sus descripciones. \\
  13.3 & Claude Code; OpenAI Codex & Contenido y recorrido INN-03 & Guillermo Castillo: Se examinó la exposición de la vida útil remanente por historia térmica, con atención a la comprensión de los supuestos y límites descritos. \\
  13.4 & Claude Code; OpenAI Codex & Contenido y recorrido INN-04 & Guillermo Castillo: Se revisó la explicación del tramo variable de la operación y la claridad de los términos relativos al costo de servir. \\
  13.5 & Claude Code; OpenAI Codex & Contenido y recorrido INN-05 & Guillermo Castillo: Se examinó la descripción de la hoja de negocio del almacenero, atendiendo a la comprensión de su propósito y funcionalidades. \\
  Anexo 13.A & Claude Code & Contratos de datos & Patricio Henríquez: Se revisó la redacción de las definiciones de los contratos de datos y la claridad de la terminología empleada. \\
  Anexo 13.B & Claude Code & Trazabilidad y alcance & Patricio Henríquez: Se examinó la presentación de la matriz de trazabilidad y la comprensión de las relaciones y del alcance descritos. \\
  Anexo 13.C & Claude Code; OpenAI Codex & Indicadores y riesgos & Patricio Henríquez: Se revisó la exposición de los indicadores y riesgos, con atención a la claridad de sus descripciones. \\
  T-19 & Claude Code; OpenAI Codex & Fichas de las cinco innovaciones & Patricio Henríquez: Se revisó la redacción, claridad y consistencia de las afirmaciones dentro de las fichas, atendiendo a la comprensión de su contenido. \\
  \midrule
  \multicolumn{4}{@{}l}{\intersemibold Subdocumento 14 \textemdash{} Ventajas, beneficios y consolidación} \\
  Introducción & Claude Code & Articulación de los documentos de la oferta & Alex Aravena y Guillermo Castillo: correspondencia con los subdocumentos y formularios citados. \\
  14.1 & Claude Code & Síntesis de la propuesta de valor y mapa de valor & Alex Aravena y Guillermo Castillo: correspondencia de cada ventaja con la decisión de diseño que la sustenta. \\
  14.2 & Claude Code & Beneficios, equilibrio, coherencia y trazabilidad & Alex Aravena y Guillermo Castillo: cotejo de las cifras con las Bases Técnicas del caso y con los subdocumentos de origen. \\
  \midrule
  \multicolumn{4}{@{}l}{\intersemibold Oferta Económica \textemdash{} Entregable 1: Propuesta económica formal} \\
  1.1 & Codex & Redacción de condiciones de precio y vigencia, y cotejo de importes, conversiones e IVA. & Maximiliano Miño. Revisó la vigencia, las conversiones, los porcentajes y el IVA contra el E-21 y el E-24. \\
  1.2 & Codex & Redacción de pagos, aceptación y liquidación variable, y cotejo del E-25 con el flujo de caja. & Tomás Pérez. Verificó hitos, calendario, montos e IVA, y cotejó la liquidación variable con el E-25 y el modelo financiero. \\
  1.3 & Codex & Redacción y síntesis del alcance, responsabilidades, beneficios y condiciones contractuales. & Maximiliano Miño. Revisó el alcance, las exclusiones, los beneficios y la coherencia de las condiciones comerciales. \\
  \midrule
  \multicolumn{4}{@{}l}{\intersemibold Oferta Económica \textemdash{} Entregable 2: Análisis económico-financiero} \\
  2.1 & Codex y Claude & Redacción y cotejo de costos, precios y márgenes con el modelo financiero & Alex Aravena: cifras contra el modelo financiero y criterios de margen \\
  2.2 & Claude & Redacción y cotejo de cifras con el modelo financiero & Tomás Pérez: cifras contra el modelo financiero y justificación de la tasa de descuento \\
  2.3 & Codex y Claude & Redacción, cotejo del flujo financiado y formato del gráfico de la Curva S & Alex Aravena: cifras contra el modelo financiero y condiciones del financiamiento \\
  2.4 & Codex y Claude & Síntesis de adquisiciones y cotejo de montos con el modelo financiero & Álvaro Catalán: montos contra el modelo financiero y coherencia con el T-11 y el T-14 \\
  2.5 & Codex y Claude & Redacción del modelo de costos operacionales y conciliación de montos & Maximiliano Miño: montos contra el modelo financiero y supuestos de servicios y gastos \\
  2.6 & Codex y Claude & Redacción del modelo de soporte y valorización de las horas de mantención & Alex Aravena: horas contra el T-14 y el T-15 y montos contra el modelo financiero \\
  2.7 & Codex y Claude & Redacción, cotejo de cifras con el modelo financiero y código de la figura de egresos & Alex Aravena: cifras contra el modelo financiero y coherencia con el T-19 \\
